\documentclass[12pt,a4paper]{report}

% ========== PACKAGES ==========
\usepackage[utf8]{inputenc}
\usepackage[T1]{fontenc}
\usepackage{times}
\usepackage[top=1in, bottom=1in, left=1.25in, right=1in]{geometry}
\usepackage{graphicx}
\usepackage{setspace}
\usepackage{titlesec}
\usepackage{tocloft}
\usepackage{fancyhdr}
\usepackage{hyperref}
\usepackage{amsmath}
\usepackage{amssymb}
\usepackage{float}
\usepackage{caption}
\usepackage{subcaption}
\usepackage{enumerate}
\usepackage{enumitem}
\usepackage{booktabs}
\usepackage{multirow}
\usepackage{array}
\usepackage{xcolor}
\usepackage{listings}

% ========== HYPERREF SETUP ==========
\hypersetup{
    colorlinks=true,
    linkcolor=black,
    filecolor=black,
    urlcolor=blue,
    citecolor=black,
    pdftitle={Project Report},
    pdfauthor={Himraj Doley, Subham Saha, Mohd Zaid, Nipujyoti Rabha},
}

% ========== PAGE STYLE ==========
\pagestyle{fancy}
\fancyhf{}
\fancyhead[L]{\leftmark}
\fancyhead[R]{\thepage}
\renewcommand{\headrulewidth}{0.4pt}

% ========== CHAPTER/SECTION FORMATTING ==========
\titleformat{\chapter}[hang]
{\normalfont\fontsize{14}{16}\bfseries\uppercase}
{\thechapter.}{1em}{}
\titlespacing*{\chapter}{0pt}{-20pt}{20pt}

\titleformat{\section}
{\normalfont\fontsize{13}{15}\bfseries}
{\thesection}{1em}{}
\titlespacing*{\section}{0pt}{12pt}{6pt}

\titleformat{\subsection}
{\normalfont\fontsize{12}{14}\bfseries}
{\thesubsection}{1em}{}
\titlespacing*{\subsection}{0pt}{10pt}{5pt}

% ========== TABLE OF CONTENTS FORMATTING ==========
\renewcommand{\cftchapfont}{\normalfont\bfseries}
\renewcommand{\cftsecfont}{\normalfont}
\renewcommand{\cftchapleader}{\cftdotfill{\cftdotsep}}
\setlength{\cftbeforechapskip}{6pt}

% ========== LINE SPACING ==========
\onehalfspacing

% ========== CODE LISTING SETUP ==========
\lstset{
    basicstyle=\ttfamily\small,
    breaklines=true,
    frame=single,
    numbers=left,
    numberstyle=\tiny,
    tabsize=4,
    showstringspaces=false
}

% ========== DOCUMENT BEGINS ==========
\begin{document}

% ========== FRONT MATTER ==========
\pagenumbering{roman}

% ========== TITLE PAGE ==========
\begin{titlepage}
    \begin{center}
        \vspace*{0.2cm}

        {\LARGE \textbf{HINDI COREFERENCE RESOLUTION}}\\[0.2cm]
        {\LARGE \textbf{FOR LOW-RESOURCE NLP USING}}\\[0.2cm]
        {\LARGE \textbf{TRANSFORMER-BASED HYBRID APPROACHES}}\\[0.7cm]
        {A Project Work Submitted in Partial Fulfillment}\\[0.2cm]
        {\large of the requirements for the Degree}\\[0.3cm]

        {\large \textit{BACHELOR OF TECHNOLOGY}}\\[0.3cm]
        {\large in}\\[0.3cm]

        {\large \textbf{COMPUTER SCIENCE AND ENGINEERING}}\\[0.5cm]

        {\large by}\\[0.3cm]

        {\large \fontsize{14}{16}\selectfont \textbf{Himraj Doley (202202023131)}}\\[0.1cm]
        {\large \fontsize{14}{16}\selectfont \textbf{Subham Saha (202202021067)}}\\[0.1cm]
        {\large \fontsize{14}{16}\selectfont \textbf{Mohd Zaid (202202023127)}}\\[0.1cm]
        {\large \fontsize{14}{16}\selectfont \textbf{Nipujyoti Rabha (202202021053)}}\\[0.4cm]

        {\large Under the supervision of}\\[0.2cm]
        {\large \fontsize{18}{16}\selectfont \textbf{Dr. Apurbalal Senapati}}\\[0.1cm]

        \includegraphics[width=2.5cm, height=2.5cm, keepaspectratio]{Picture1.png}\\[0.3cm]

        {\large \textbf{DEPARTMENT OF COMPUTER SCIENCE \& ENGINEERING}}\\[0.1cm]
        {\large \textbf{CENTRAL INSTITUTE OF TECHNOLOGY KOKRAJHAR}}\\[0.1cm]

        \includegraphics[width=8cm, height=0.7cm, keepaspectratio]{55.png}\\[0.1cm]

        {\fontsize{10}{10}\selectfont \textbf{(A Centrally Funded Institute under Ministry of HRD, Govt. of India)}}\\[0.1cm]
        {\fontsize{10}{10}\selectfont \textbf{BODOLAND TERRITORIAL AREAS DISTRICTS :: KOKRAJHAR :: ASSAM :: 783370}}\\[0.05cm]
        {\fontsize{10}{10}\selectfont \textbf{Website: www.cit.ac.in}}\\[0.05cm]
        {\fontsize{10}{10}\selectfont \textbf{BTAD, ASSAM}}\\[0.1cm]
        {\large \textbf{MAY 2026}}
    \end{center}
\end{titlepage}


% ========== CERTIFICATE OF APPROVAL PAGE ==========
\newpage
\thispagestyle{empty}
\begin{center}
    \includegraphics[width=3.2cm, height=3.2cm, keepaspectratio]{Picture1.png}\\[0.2cm]
    {\large \textbf{DEPARTMENT OF COMPUTER SCIENCE AND ENGINEERING}}\\[0.2cm]
    \includegraphics[width=8cm, height=0.7cm, keepaspectratio]{55.png}\\[0.2cm]
    {\large \textbf{CENTRAL INSTITUTE OF TECHNOLOGY KOKRAJHAR}}\\[0.2cm]
    {\large (An Autonomous Institute under MHRD) Kokrajhar -- 783370, BTAD, Assam, India}\\[1.5cm]

    {\Large \textbf{\underline{CERTIFICATE OF APPROVAL}}}\\[1cm]
\end{center}

\noindent This is to certify that the work presented in this project report, titled \textbf{``Hindi Coreference Resolution for Low-Resource NLP using Transformer-Based Hybrid Approaches''} and submitted by \textbf{Himraj Doley, Subham Saha, Mohd Zaid, and Nipujyoti Rabha} to the Department of Computer Science \& Engineering, has been carried out under our direct supervision and guidance throughout its entirety.\\[0.5cm]

\noindent The project has been developed in accordance with the academic regulations of \textbf{Central Institute of Technology}, and we recommend that this work be accepted in partial fulfillment of the requirements for the degree of Bachelor of Technology in Computer Science and Engineering.\\[2cm]

\noindent \textbf{Supervisor}\\[0.3cm]
\textbf{Dr. Apurbalal Senapati}\\
(HOD, Dept. of CSE)\\[1cm]

\noindent Date: \rule{3cm}{0.4pt}\\[0.3cm]
Place: Kokrajhar, BTAD, Assam

% ========== CERTIFICATE OF BOARD OF EXAMINERS PAGE ==========
\newpage
\thispagestyle{empty}
\begin{center}
    \includegraphics[width=3.2cm, height=3.2cm, keepaspectratio]{Picture1.png}\\[0.2cm]
    {\large \textbf{DEPARTMENT OF COMPUTER SCIENCE AND ENGINEERING}}\\[0.2cm]
    \includegraphics[width=8cm, height=0.7cm, keepaspectratio]{55.png}\\[0.2cm]
    {\large \textbf{CENTRAL INSTITUTE OF TECHNOLOGY KOKRAJHAR}}\\[0.2cm]
    {\large (An Autonomous Institute under MHRD) Kokrajhar -- 783370, BTAD, Assam, India}\\[1cm]

    {\Large \textbf{\underline{Certificate by the Board of Examiners}}}\\[0.8cm]
\end{center}

\noindent This is to certify that the project report entitled \textbf{``Hindi Coreference Resolution for Low-Resource NLP using Transformer-Based Hybrid Approaches''}, submitted by \textbf{Himraj Doley, Subham Saha, Mohd Zaid, and Nipujyoti Rabha} to the Department of Computer Science and Engineering, Central Institute of Technology Kokrajhar, has been duly examined and evaluated by the Board of Examiners.\\[0.2cm]

\noindent The project has been conducted and documented in conformity with the academic regulations of \textbf{Central Institute of Technology}, and qualifies for acceptance in partial fulfillment of the requirements for the degree of Bachelor of Technology in \textbf{Computer Science and Engineering}.\\[0.7cm]

\noindent \textbf{Project Co-ordinator}\\[0.5cm]
\rule{4cm}{0.4pt}\\[0.1cm]
\textbf{Mr. Mithun Karmakar}\\
(Assistant Professor, Dept. of CSE)\\[0.3cm]

\noindent \textbf{Countersigned by}\\[0.1cm]

\noindent \textbf{Head Of Department}\\[0.5cm]
\rule{4cm}{0.4pt}\\[0.1cm]
\textbf{Dr. Apurbalal Senapati}\\
(HOD, Dept. of CSE)


% ========== ACKNOWLEDGEMENT ==========
\newpage
\chapter*{\fontsize{20}{24}\selectfont ACKNOWLEDGEMENT}
\addcontentsline{toc}{chapter}{ACKNOWLEDGEMENT}

\noindent Completing this project has been one of the most challenging and rewarding experiences of our undergraduate journey, and it would not have been possible without the support and encouragement of many people along the way.

\noindent We owe our deepest gratitude to our project supervisor, \textbf{Dr. Apurbalal Senapati}, whose guidance was the cornerstone of everything we accomplished here. His patience during the difficult early stages --- particularly when our first approach hit a fundamental dead end --- and his willingness to help us think through a more principled alternative made all the difference. The decision to pivot from an unsupervised embedding-based method to a supervised hybrid architecture came directly out of those conversations, and the quality of the final system reflects his mentorship at every step.\\[0.5cm]

\noindent As Head of the Department of Computer Science and Engineering at \textbf{Central Institute of Technology Kokrajhar}, \textbf{Dr. Apurbalal Senapati} also ensured that we had access to the resources and academic environment necessary to pursue research of this scope. We are likewise grateful to all the faculty members of the department for their feedback during project reviews and for being genuinely approachable whenever we needed guidance outside scheduled hours.\\[0.5cm]

\noindent A project of this nature demands sustained teamwork, and we want to acknowledge the contribution of every member of our group. The manual annotation effort alone --- going through Hindi documents one by one, marking up coreference spans, and building a dataset of 76 documents and 548 mention annotations from scratch --- required hours of careful work spread over several weeks. That effort would simply not have come together without everyone pulling their weight consistently.\\[0.5cm]

\noindent We are also thankful to our families and friends, whose quiet encouragement kept us going through the long stretches when progress felt slow. Finally, we acknowledge with gratitude the open-source community behind HuggingFace Transformers, the AI4Bharat team for making IndicBERTv2 publicly available, and the broader Hindi NLP research community whose prior work gave us the foundation to build upon.

\vspace{1cm}
\begin{flushright}
\textbf{Submitted By:}\\
Himraj Doley    (202202023131)\\
Subham Saha     (202202021067)\\
Mohd Zaid       (202202023127)\\
Nipujyoti Rabha (202202021053)
\end{flushright}


% ========== ABSTRACT ==========
\newpage
\chapter*{\fontsize{16}{19}\selectfont ABSTRACT}
\addcontentsline{toc}{chapter}{ABSTRACT}

\noindent This report documents the development of a \textbf{Hindi Coreference Resolution system} designed to operate under low-resource conditions. The system combines contextual embeddings from a pretrained Indic language model with carefully engineered Hindi-specific linguistic features, forming a hybrid architecture suited to a setting where annotated training data is scarce.

\noindent The project unfolded in two distinct phases, each of which contributed meaningfully to the final outcome. In the \textbf{first phase}, we trained a compact BERT-style transformer from scratch using the Masked Language Modeling (MLM) objective on a corpus of approximately 172,000 Hindi Wikipedia documents sourced from Kaggle. A SentencePiece Byte Pair Encoding tokenizer with a vocabulary size of 32,000 was trained on the same data. Although the resulting embeddings exhibited reasonably coherent semantic structure in neighbourhood analyses, the approach broke down entirely when applied to coreference resolution. Pronouns such as \textit{vah} (he/she), \textit{usne} (he/she did), and \textit{vahaan} (there) were almost never linked to their correct antecedents, regardless of the similarity threshold used for clustering. The core problem was conceptual: coreference resolution is fundamentally a referential task that cannot be reduced to semantic proximity. Recognising this, we formally abandoned Phase~1 and treated its failure as a research finding in its own right.

\noindent The \textbf{second phase} adopted a supervised, hybrid architecture. Rather than training a language model from scratch, we built upon \textbf{AI4Bharat's IndicBERTv2 (MLM-only)}, which already encodes substantial Hindi linguistic knowledge through large-scale pretraining on Indic corpora. Because no Hindi coreference dataset of the required format was publicly available, we manually created one: a corpus of \textbf{76 documents} carrying \textbf{548 mention annotations}, each labelled with a cluster identifier, mention span boundaries, and mention type. Using this dataset, we trained a \textbf{pairwise mention classifier} built on top of IndicBERTv2 contextual representations, augmented by handcrafted features capturing Hindi-specific properties such as pronoun type, plural pronoun status, object-case markers, clause-level heuristics, and mention compatibility signals. The classifier was a \textbf{Logistic Regression} model trained with balanced class weights, evaluated through five-fold stratified cross-validation and achieving approximately \textbf{82\% overall pairwise accuracy}. Predicted coreferent mention pairs were then merged into entity clusters using a \textbf{Union-Find} algorithm. The complete pipeline was evaluated using both pairwise classification metrics (precision, recall, F1) and the MUC coreference scoring scheme.

\noindent Taken together, this project demonstrates that a viable coreference pipeline for a low-resource language can be constructed through a combination of informed model selection, deliberate dataset creation, and domain-aware feature engineering. The report includes an honest account of the limitations of the current system --- including its dependence on pre-supplied mention spans, the ceiling imposed by logistic regression, and the small size of the annotated corpus --- along with a concrete roadmap for future development.

\vspace{0.5cm}
\noindent \textbf{Keywords:} Coreference Resolution, Hindi NLP, IndicBERT, IndicBERTv2, Masked Language Modeling, Low-Resource NLP, Pairwise Classification, Logistic Regression, Union-Find Clustering, Transformer Models, Manual Annotation, Linguistic Feature Engineering

% ========== TABLE OF CONTENTS ==========
\newpage
\tableofcontents

% ========== LIST OF FIGURES ==========
\newpage
\listoffigures
\addcontentsline{toc}{chapter}{LIST OF FIGURES}

% ========== LIST OF TABLES ==========
\newpage
\listoftables
\addcontentsline{toc}{chapter}{LIST OF TABLES}

% ========== MAIN CONTENT ==========
\newpage
\pagenumbering{arabic}
\setcounter{page}{1}


% ========== CHAPTER 1: INTRODUCTION ==========
\chapter{INTRODUCTION}

\section{Background}

Over the past decade, Natural Language Processing has undergone a profound shift. The arrival of large pretrained transformer models --- BERT \cite{devlin2018bert}, RoBERTa \cite{liu2019roberta}, and their numerous successors --- fundamentally changed what was possible across almost every NLP benchmark. Yet this progress has not been evenly distributed. The bulk of research investment and annotated data continues to concentrate on English, with a modest extension to a few high-resource European languages. For Hindi, a language spoken natively by over 530 million people and understood by perhaps a hundred million more as a second language --- making it one of the most widely spoken languages on earth --- the NLP landscape remains considerably underdeveloped, particularly for tasks that require discourse-level understanding.

\textbf{Coreference resolution} is one such task. At its core, coreference resolution asks a deceptively simple question: which expressions in a text are talking about the same thing? The goal is to identify all referring expressions --- pronouns, noun phrases, named entities --- that point to the same real-world entity and to group them into what are called \textit{coreference chains} or entity clusters. A brief example makes this concrete:

\begin{center}
\textit{``Ram school gaya. Vah khush tha.''}\\
(Ram went to school. He was happy.)
\end{center}

\noindent A coreference resolver working on this passage must determine that \textit{Ram} and \textit{Vah} (he) both refer to the same individual and therefore belong in the same cluster. What appears trivial to a human reader is actually a non-trivial inference for a computational system: it must understand pronoun semantics, discourse continuity, and the morphological properties of a richly inflected language, all simultaneously. The ability to perform this kind of entity tracking is a prerequisite for many downstream applications including machine translation, question answering, dialogue systems, information extraction, and document summarisation. A system that cannot correctly resolve coreferences will introduce systematic errors in all of these tasks.

\section{Why Hindi Is Particularly Challenging}

While coreference resolution is hard in general, Hindi introduces a distinct set of complications that make the problem considerably more difficult than it is for English:

\begin{itemize}
    \item \textbf{Morphological richness:} Hindi nouns, pronouns, and verbs carry inflections for gender, number, case, and tense. The pronoun \textit{vah}, for instance, can refer to a person of any gender, which means the system cannot resolve it on morphological grounds alone and must rely on broader contextual cues.

    \item \textbf{Pro-drop behaviour:} Hindi frequently omits the subject pronoun when it can be inferred from the verb or discourse context. This means a coreference system must track implicit arguments that are never explicitly stated in the text --- a capability that purely surface-level approaches cannot support.

    \item \textbf{Relatively free word order:} Although Subject-Object-Verb is the canonical order in Hindi, deviations are common in spoken and written registers, making it harder to rely on positional heuristics for identifying antecedents.

    \item \textbf{Devanagari script complexity:} Conjunct characters, matras, and zero-width joiners mean that even tokenisation is non-trivial. A tokeniser trained on Latin-script data is likely to mishandle Devanagari in ways that propagate errors downstream.

    \item \textbf{Severe shortage of annotated data:} This is arguably the biggest obstacle. English benefits from resources like the OntoNotes 5.0 corpus, which contains tens of thousands of coreference-annotated documents spread across multiple genres. Hindi has no comparable resource. The few annotated corpora that exist in Hindi are small, domain-restricted, and not always publicly accessible.
\end{itemize}

These challenges interact with each other in ways that make off-the-shelf solutions difficult to apply. A system designed for English, even a powerful multilingual one, will often struggle with Hindi precisely because it was never exposed to the morphological and discourse patterns that Hindi coreference demands.

\section{How This Project Came Together}

This project did not follow a straight line from problem statement to working system. It evolved through two distinct phases, and being honest about that progression --- including a significant failure --- is central to what this report tries to convey.

\subsection{Phase 1: Training a BERT Model from Scratch}

Our initial hypothesis was that training a BERT-style transformer directly on Hindi Wikipedia data would give us contextual representations rich enough to support coreference resolution through unsupervised similarity clustering. We assembled a corpus of roughly 172,000 Hindi Wikipedia documents from Kaggle, trained a SentencePiece BPE tokenizer with a 32,000-word vocabulary, and pretrained a compact four-layer transformer using the Masked Language Modeling objective.

The embeddings that came out of this process were not without value --- semantically related words did cluster together to a reasonable degree in visualisations. But when we tried to use cosine similarity between mention embeddings as a basis for coreference linking, the system broke down completely. Pronouns like \textit{vah}, \textit{usne}, and \textit{vahaan} were consistently unlinked from their antecedents, regardless of how we set the similarity threshold. The problem was not a tuning issue. It reflected a more fundamental misunderstanding on our part: coreference is not about semantic similarity. Two expressions can be completely semantically unlike --- \textit{Ram} and \textit{vah} share no meaning overlap --- and yet be coreferent. The relationship between them is referential, not semantic, and no amount of embedding geometry can bridge that gap without additional structural information. We concluded Phase~1 with a clear negative result and a much better understanding of what the task actually requires.

\subsection{Phase 2: A Supervised Hybrid Architecture}

Armed with that understanding, we redesigned the system from the ground up. Instead of training our own language model, we adopted \textbf{IndicBERTv2} from the AI4Bharat project, which had already been pretrained on large, high-quality Indic language corpora and handled Devanagari tokenisation well. This gave us strong contextual representations without requiring us to reproduce years of large-scale pretraining on limited hardware.

The remaining challenge was the absence of labelled training data. We addressed this by manually annotating a dataset of 76 Hindi documents with 548 coreference mention spans, each assigned a cluster identifier and mention type. Using this dataset, we formulated coreference as a \textbf{pairwise classification problem}: for every pair of mentions within a document, a classifier predicts whether they refer to the same entity. The feature vector fed to the classifier combined IndicBERTv2 contextual embeddings with a set of handcrafted features encoding Hindi-specific pronoun behaviour, number agreement signals, clause-level heuristics, and mention distance. A Logistic Regression classifier with balanced class weights was trained on these features and evaluated through five-fold cross-validation. Predicted coreferent pairs were finally assembled into entity clusters using a Union-Find algorithm.

\section{Contributions of This Work}

The specific contributions that this project makes are as follows:

\begin{itemize}
    \item A documented, principled analysis of why embedding-based unsupervised clustering fails for Hindi coreference, serving as a useful negative result for others working in this space.
    \item A manually annotated Hindi coreference dataset of 76 documents and 548 mention spans, created from scratch to address the resource gap.
    \item A working end-to-end coreference pipeline for Hindi that integrates IndicBERTv2 contextual representations, domain-specific linguistic features, pairwise logistic regression classification, and Union-Find cluster reconstruction.
    \item Empirical evidence, through an ablation study, that handcrafted Hindi linguistic features add approximately 11 percentage points of classification accuracy over embedding similarity alone.
    \item An honest assessment of where the system falls short and a concrete set of directions for follow-on work.
\end{itemize}

\section{Organisation of the Report}

Chapter~2 sets out the motivation and objectives of the project in detail. Chapter~3 surveys the relevant prior work across transformer models, Indic language NLP, and coreference resolution systems. Chapter~4 describes the complete methodology for both phases. Chapter~5 presents the experimental results and their interpretation. Chapter~6 offers concluding remarks along with a frank discussion of the system's limitations. Chapter~7 outlines directions for future research.


% ========== CHAPTER 2: MOTIVATION & OBJECTIVE ==========
\chapter{MOTIVATION \& OBJECTIVE}

\section{Motivation}

\subsection{The Resource Gap in Hindi NLP}

Hindi occupies an unusual position in the global language landscape. It is the official language of India, spoken natively by over 530 million people and used as a second language by roughly another 100 million. By sheer speaker count, it ranks among the top three most spoken languages in the world. And yet, within the field of Natural Language Processing, Hindi is consistently treated as a low-resource language --- not because it is rare or obscure, but because the research community has not produced annotated corpora for it at anything close to the scale that exists for English.

The contrast becomes stark when one looks at specific tasks. For English coreference resolution, the OntoNotes 5.0 corpus provides thousands of carefully annotated documents spanning multiple genres, and researchers have had access to it for over a decade. For Hindi, there is no equivalent. The handful of annotated coreference resources that do exist are small, narrowly scoped, and largely inaccessible outside academia. This gap directly limits what NLP systems built for Hindi can accomplish: machine translation systems produce incorrect pronoun mappings because they cannot track entity gender; question-answering systems fail on multi-sentence passages because they cannot link anaphors to their referents; summarisation systems produce incoherent output because they lose track of entities across paragraph boundaries.

Our motivation for this project grew out of a genuine interest in closing some part of this gap, even if incrementally. We wanted to understand what could be achieved for Hindi coreference with limited supervision, restricted compute, and no access to large proprietary datasets --- the conditions that most researchers in low-resource settings actually face.

\subsection{Why Coreference Resolution Matters}

Coreference resolution is sometimes described as a subtask within NLP, but that framing understates its importance. It is better understood as a foundational capability that other language understanding tasks depend on. Several examples illustrate this:

\begin{itemize}
    \item \textbf{Machine Translation:} Hindi and English differ substantially in how gender is expressed. Correctly translating a Hindi pronoun into English requires knowing what entity it refers to, because the antecedent's grammatical properties determine the correct English pronoun. Without coreference information, translation systems are forced to guess, and they frequently guess wrong.

    \item \textbf{Question Answering:} A system asked ``Who went to the market?'' after reading a paragraph must be able to trace the pronoun through the text back to the named entity it refers to. Multi-hop reasoning of this kind is impossible without entity tracking.

    \item \textbf{Document Summarisation:} A good summary must refer to entities consistently and coherently. Without knowing which pronouns and noun phrases refer to the same person or object, an automatic summariser either produces dangling references or unnecessarily repeats full names.

    \item \textbf{Dialogue Systems and Virtual Assistants:} When a user says ``Tell me more about her'' in a conversation, the system needs to resolve ``her'' against the dialogue history. This is a real-time coreference resolution problem, and it becomes particularly challenging in Hindi where pronouns are underspecified for gender.

    \item \textbf{Information Extraction and Knowledge Graphs:} Building a structured knowledge base from Hindi text requires consolidating multiple mentions of the same entity across a document or corpus. Coreference resolution provides exactly this consolidation.
\end{itemize}

In each of these cases, a system without coreference resolution capability hits a hard performance ceiling that cannot be overcome by improving other components in isolation.

\subsection{The Inadequacy of Generic Multilingual Models for This Task}

One might reasonably ask whether an existing multilingual model such as mBERT or XLM-RoBERTa could simply be applied to Hindi coreference. The short answer is: not without fine-tuning, and fine-tuning requires labelled data that does not currently exist at scale. Beyond the data problem, generic multilingual models carry additional limitations for Hindi specifically:

\begin{itemize}
    \item Their vocabulary and tokenisation are shared across many languages, which means Devanagari text is often over-segmented into subword units that do not align cleanly with morphological boundaries.
    \item Training signal for Hindi is diluted across the many languages these models cover, which reduces their sensitivity to Hindi-specific morphological and discourse patterns.
    \item None of these models have been fine-tuned for Hindi coreference, meaning their representations, however strong in general, are not oriented toward the referential relationships that coreference requires.
\end{itemize}

This project was specifically motivated by the question of whether a system built around a Hindi-focused pretrained model, combined with explicit linguistic feature engineering, could do better than the generic multilingual approach in this severely data-limited setting.

\subsection{The Academic Value of the Failure Documentation}

There is also a less obvious dimension to our motivation. When our first approach --- training a BERT model from scratch and using its embeddings for unsupervised coreference clustering --- failed, we initially regarded this as a setback. On reflection, however, the failure had genuine analytical value. It produced a concrete, well-understood negative result: embedding-based semantic similarity is not sufficient for coreference resolution, and a small scratch-trained model on limited data is not sufficient for embedding quality either. Documenting this finding carefully, with a root-cause analysis, serves the research community: it saves future students and researchers from repeating the same experiment under the mistaken belief that semantic similarity and referential identity are related problems.

\section{Objectives}

\subsection{Primary Objectives}

\begin{enumerate}
    \item \textbf{Empirically evaluate the limits of unsupervised coreference:} Conduct a systematic investigation into whether a scratch-trained BERT model, combined with embedding similarity clustering, can resolve coreferences in Hindi. Document the outcome and its causes clearly.

    \item \textbf{Build a functional end-to-end Hindi coreference pipeline:} Develop a complete system that takes Hindi text with annotated mention spans as input and produces entity clusters as output, with all intermediate steps documented and reproducible.

    \item \textbf{Address the data scarcity problem through manual annotation:} Create a Hindi coreference dataset from scratch, annotating 76 documents with 548 mention spans, cluster identifiers, and mention types, and make this dataset available for reproducible evaluation.

    \item \textbf{Design a hybrid architecture that combines neural and linguistic knowledge:} Exploit the contextual representation quality of IndicBERTv2 while supplementing it with handcrafted features that encode Hindi grammar --- a combination motivated by the literature on low-resource NLP.
\end{enumerate}

\subsection{Technical Objectives}

\begin{enumerate}
    \item Train a SentencePiece BPE tokenizer with a vocabulary of 32,000 subword units on the Hindi Wikipedia corpus of approximately 172,000 documents.
    \item Pretrain a compact BERT-style transformer using the MLM objective on this corpus and evaluate its embedding quality for coreference clustering (Phase~1).
    \item Integrate IndicBERTv2 as the contextual encoder for mention span representations in Phase~2, using mean pooling over span token hidden states.
    \item Construct a feature vector for each mention pair that captures both embedding-level similarity and Hindi-specific linguistic signals including pronoun detection, plurality markers, object-case flags, clause-level co-occurrence, and mention compatibility.
    \item Train and evaluate a Logistic Regression pairwise classifier on the manually annotated dataset using five-fold stratified cross-validation.
    \item Implement Union-Find cluster reconstruction from predicted coreferent mention pairs.
    \item Report system performance using standard pairwise classification metrics and the MUC coreference scoring scheme.
\end{enumerate}

\subsection{Research Objectives}

\begin{enumerate}
    \item Understand, at a mechanistic level, why semantic similarity is an insufficient basis for coreference resolution, and articulate this finding in terms that generalise beyond Hindi.
    \item Quantify the individual contribution of the linguistic feature group relative to embedding features alone through an ablation study.
    \item Identify the primary bottlenecks in the Phase~2 system --- whether they lie in dataset size, mention detection quality, classifier capacity, or cluster reconstruction logic --- and characterise them honestly.
    \item Situate this work within the broader literature on low-resource NLP and hybrid neural-symbolic system design, identifying clearly where our approach aligns with and diverges from established practice.
\end{enumerate}


% ========== CHAPTER 3: RELATED WORK ==========
\chapter{RELATED WORK}

This chapter reviews the prior work most directly relevant to our project, organised across four areas: transformer-based language models, Indic language modelling, coreference resolution systems, and strategies for building NLP systems under low-resource conditions.

\section{Transformer-Based Language Models}

\subsection{BERT and the Pretrain-Then-Finetune Paradigm}

The publication of BERT by Devlin et al.\ \cite{devlin2018bert} marked a turning point in how the NLP community approached representation learning. The key insight was that training a deep bidirectional transformer using Masked Language Modeling (MLM) and Next Sentence Prediction (NSP) on large unlabelled text corpora produces representations that transfer extremely well to downstream tasks through fine-tuning. Prior to BERT, the dominant paradigm involved training task-specific architectures from scratch on relatively small labelled datasets; BERT demonstrated that a single pretrained model, adjusted with a small amount of task-specific supervision, could outperform specialised systems across a wide range of benchmarks.

The bidirectionality of BERT is particularly important for coreference resolution. An antecedent may appear either before or after the expression it anchors --- a pronoun might follow its referent by several sentences, or in less common constructions, precede it. A left-to-right language model fundamentally cannot capture context from both directions simultaneously, whereas the masked language modelling objective forces each token to be predicted using the full surrounding context. This property makes BERT-style representations a natural choice for any discourse-level task.

Liu et al.\ \cite{liu2019roberta} subsequently investigated the sensitivity of BERT's performance to training decisions through their RoBERTa model. They found that removing the NSP objective, training on significantly larger batches, and applying dynamic masking during training all led to consistent improvements. Their findings were relevant to our Phase~1 work: the fact that even a carefully designed large-scale BERT variant benefits from substantial data reinforced our expectation that a small scratch-trained model on 172,000 documents would produce embeddings of limited quality for a complex discourse task.

\subsection{Why Small Scratch-Trained Models Fall Short}

Research on transfer learning in low-resource settings has consistently shown that transformer models pretrained on small corpora produce representations that are considerably weaker than their large-scale counterparts, particularly for tasks that require deep linguistic or discourse-level understanding \cite{lauscher2020zero}. The reason is structural: the value of large-scale pretraining lies not merely in learning word statistics but in acquiring a broadly useful inductive bias about language that can be efficiently specialised with task-specific data. A model trained on a limited corpus does not acquire this generalised bias to the same degree, which means it requires proportionally more supervised data to perform well on challenging tasks. For our Phase~1 model, this translated directly into poor embedding quality for pronoun-antecedent linking, a relationship that requires sensitivity to discourse structure that the model simply had not learned.

\section{Indic Language Models}

\subsection{IndicBERT and the IndicNLPSuite}

Kakwani et al.\ \cite{kakwani2020indicnlpsuite} presented one of the first systematic efforts to build NLP infrastructure for Indian languages at scale. The IndicNLPSuite comprised a large monolingual corpus (IndicCorp) spanning 12 Indian languages, a set of evaluation benchmarks for tasks including text classification, named entity recognition, and natural language inference, and IndicBERT --- a multilingual BERT model pretrained specifically on Indian language text. IndicBERT represented a significant step forward over mBERT for Indic language tasks, primarily because its training data was not diluted across unrelated language families. Nevertheless, limitations in tokenisation quality and corpus size meant that performance on morphologically demanding tasks remained inconsistent.

\subsection{IndicBERTv2}

The same team at AI4Bharat subsequently released IndicBERTv2 \cite{doddapaneni2023indicbert}, a substantially improved model pretrained on a larger and more carefully curated version of the IndicCorp corpus (IndicCorp v2). The MLM-only variant of IndicBERTv2 --- which drops the NSP objective in line with insights from RoBERTa --- showed improved performance across multiple benchmarks, particularly on tasks involving morphological sensitivity. For our purposes, two properties of IndicBERTv2 were especially relevant. First, its Devanagari tokenisation was noticeably more robust, producing subword units that aligned more cleanly with Hindi morphological boundaries. Second, its pretraining on a large Hindi corpus meant that the contextual representations it produces for Hindi text carry significantly more linguistic information than anything our small scratch-trained model could generate. These properties made IndicBERTv2 the natural choice for the contextual encoder in our Phase~2 system.

\subsection{MuRIL}

Khanuja et al.\ \cite{khanuja2021muril} released MuRIL (Multilingual Representations for Indian Languages), a model pretrained on 17 Indian languages together with their transliterated and machine-translated counterparts. MuRIL outperformed mBERT on several Indic NLP benchmarks by a considerable margin, largely because its training data was richer and more domain-diverse. However, MuRIL's larger parameter count and dependence on transliteration pipelines as part of its training setup made it less practical for our constrained hardware environment. We considered it as an alternative to IndicBERTv2 but ultimately stayed with the latter given its smaller footprint and its direct optimisation for the MLM task on Indic text without transliteration augmentation.

\section{Coreference Resolution}

\subsection{Early Rule-Based and Statistical Approaches}

The problem of pronoun and noun phrase resolution has been studied in computational linguistics since at least the early 1990s. Lappin and Leass \cite{lappin1994algorithm} described one of the early influential algorithms for pronominal anaphora resolution, relying on syntactic structure and salience factors such as sentence recency and grammatical role. While limited in scope by modern standards, their work established the key intuition that pronouns tend to resolve to the most salient entities in the preceding discourse context --- an insight that still informs feature engineering decisions today, including in our own system.

Soon et al.\ \cite{soon2001machine} formalised the pairwise mention model, which remains the conceptual backbone of our Phase~2 classifier. Their approach cast coreference as a binary classification problem: for each anaphor, a classifier examines preceding mention candidates and predicts which, if any, is coreferent with it. The features they used --- string match, pronoun detection, distance, definiteness, grammatical role --- were relatively simple by current standards, but the architectural choice to decompose coreference into pairwise decisions proved robust and has influenced system designs ever since. Our feature set is directly inspired by this tradition, extended with Hindi-specific morphological and discourse signals.

\subsection{Neural Coreference Systems}

The transition to neural approaches began gradually, with early systems replacing hand-engineered features with learned representations while retaining the pairwise mention-scoring architecture. The most significant leap came with Lee et al.\ \cite{lee2017end}, who introduced the first end-to-end neural coreference system. Rather than working from pre-detected mention spans, their system treated all text spans up to a maximum length as candidate mentions and jointly trained a mention scoring network and an antecedent scoring network using a single objective. This formulation eliminated the cascade of errors that plagued pipeline systems, where mention detection mistakes would propagate to the coreference scoring stage.

Joshi et al.\ \cite{joshi2019bert} subsequently replaced the custom span representations in Lee et al.'s system with BERT contextual embeddings, producing a dramatic performance improvement on the OntoNotes English benchmark and establishing the current state of the art for English coreference. For English, systems of this type now achieve F1 scores above 80\% on OntoNotes. The gap between these figures and what is currently achievable for Hindi is a direct consequence of the resource disparity: the OntoNotes training set for English contains tens of thousands of annotated documents, while our Hindi dataset contains 76.

\subsection{Unsupervised Clustering Approaches and Their Limitations}

Several researchers have explored clustering-based coreference approaches that avoid the need for labelled data. The general strategy is to compute contextual embeddings for all candidate mentions and group them using distance metrics in embedding space. While these approaches have met with some success in very specific settings --- particularly when mentions are consistently described using similar lexical forms --- they share a structural weakness that our Phase~1 experience confirmed directly. Embedding similarity captures semantic relatedness, not referential identity. A pronoun and its antecedent often have no semantic similarity whatsoever; what links them is discourse structure, not meaning. Unsupervised clustering therefore tends to group semantically similar non-coreferent mentions together while failing to link dissimilar coreferent ones. No threshold choice resolves this tension, because the problem is with the underlying signal, not its calibration.

\subsection{Hindi-Specific Coreference Work}

Research targeting Hindi coreference specifically is sparse. Mujadia et al.\ \cite{mujadia2016hindi} provided annotation guidelines and a small CoNLL-format Hindi coreference dataset, which remains one of the few documented resources in this space. Earlier, Dakwale et al.\ \cite{dakwale2013anaphora} described a hybrid anaphora resolution system for Hindi that combined syntactic rules with statistical decision components, reporting moderate precision on their test set. Neither of these systems has been followed up with large-scale neural approaches, primarily because the annotated data needed to train such systems does not exist. Our project works in the same research space, building on these earlier contributions while taking a different hybrid approach that leverages pretrained transformer representations unavailable at the time of those publications.

\section{Low-Resource NLP Strategies}

A review of the survey literature on low-resource NLP \cite{hedderich2021survey} reveals a set of strategies that have proven consistently useful when labelled data is scarce. Four of these are directly instantiated in our Phase~2 system:

\begin{itemize}
    \item \textbf{Transfer learning from pretrained models:} Rather than training from scratch, leveraging a language-specific pretrained model (IndicBERTv2) provides strong generic linguistic representations that small task-specific datasets can specialise without overfitting. This is the single most impactful design decision in our Phase~2 architecture.

    \item \textbf{Feature engineering as a form of data efficiency:} When labelled examples are too few to teach a model the relevant patterns from scratch, hand-engineered features can encode prior linguistic knowledge directly. In our case, features such as the pronoun flag and plurality signal give the classifier information about Hindi grammar that it would need thousands of examples to learn from data alone.

    \item \textbf{Simple classifiers with informative features:} In genuinely low-data regimes, high-capacity neural classifiers tend to overfit. Logistic regression and support vector machines, by contrast, have lower variance and are more robust to small sample sizes, provided the feature representation is informative \cite{lappin1994algorithm}. This is the rationale behind our choice of logistic regression over a neural pairwise scorer.

    \item \textbf{Manual corpus creation:} When no suitable dataset exists, building one --- even at small scale --- is a legitimate and necessary research contribution. The 76-document annotated corpus we created is small, but it enabled supervised training and provided a concrete evaluation benchmark that did not previously exist.
\end{itemize}

All four of these strategies are grounded in the established literature, and their application in our system reflects a deliberate alignment with best practices for the low-resource setting rather than arbitrary design choices.


% ========== CHAPTER 4: PROPOSED METHODOLOGY ==========
\chapter{PROPOSED ALGORITHM / DESIGN / METHODOLOGY}

This chapter describes the complete methodology underlying both phases of the project. Phase~1 covers the scratch-BERT pretraining and unsupervised clustering experiment, including its design rationale and the analysis of its failure. Phase~2 covers the supervised hybrid architecture that replaced it, describing each component of the pipeline in detail.

\section{Phase 1: Scratch BERT Pretraining and Unsupervised Clustering}

\subsection{Corpus and Preprocessing}

The training data for Phase~1 was drawn from a Kaggle dataset of Hindi Wikipedia documents comprising approximately \textbf{172,000 files}. The corpus was used in its raw, unannotated form, since the Phase~1 objective was unsupervised pretraining. Before tokenisation and training, a cleaning pipeline was applied to each document:

\begin{itemize}
    \item Wikipedia markup syntax, hyperlinks, and residual HTML tags were stripped using regular expressions.
    \item Devanagari Unicode characters were normalised to a consistent canonical form to eliminate encoding inconsistencies.
    \item Sentence boundaries were identified primarily using the Devanagari \textit{daṇḍa} (``।'') as the sentence terminal marker, supplemented by punctuation heuristics for sentences ending with question marks or exclamation marks.
    \item Documents with fewer than ten tokens after cleaning, and those containing a majority of non-Hindi characters, were discarded to avoid introducing noise into the training data.
\end{itemize}

\subsection{SentencePiece BPE Tokenizer}

Rather than using a pre-existing tokenizer, we trained a \textbf{SentencePiece} tokenizer using Byte Pair Encoding (BPE) directly on the cleaned Hindi corpus. The configuration used is summarised in Table~\ref{tab:sp_config}.

\begin{table}[H]
\centering
\caption{SentencePiece Tokenizer Configuration}
\label{tab:sp_config}
\begin{tabular}{|l|l|}
\hline
\textbf{Parameter} & \textbf{Value} \\
\hline
Algorithm & BPE (Byte Pair Encoding) \\
Vocabulary size & 32,000 \\
Character coverage & 0.9995 \\
Special tokens & \texttt{[PAD]}, \texttt{[UNK]}, \texttt{[CLS]}, \texttt{[SEP]}, \texttt{[MASK]} \\
Training corpus & Hindi Wikipedia ($\sim$172K documents) \\
\hline
\end{tabular}
\end{table}

\noindent SentencePiece was preferred over WordPiece for a specific reason: it operates directly on raw Unicode text without requiring whitespace-based pre-tokenisation. This matters for Devanagari because whitespace boundaries in Hindi text do not always correspond to meaningful morphological units --- certain postpositions and conjunctions attach to preceding words without intervening spaces --- and WordPiece's assumption that space marks a word boundary can produce poor segmentations as a result.

\subsection{BERT-Style Transformer Architecture}

The model architecture was a standard BERT encoder configuration, instantiated using the \texttt{BertForMaskedLM} class from the HuggingFace Transformers library \cite{wolf2020transformers}. Given the compute constraints of the Kaggle and Google Colab environments available to us, a deliberately compact architecture was chosen, as detailed in Table~\ref{tab:bert_config}.

\begin{table}[H]
\centering
\caption{Scratch BERT Model Configuration}
\label{tab:bert_config}
\begin{tabular}{|l|l|}
\hline
\textbf{Parameter} & \textbf{Value} \\
\hline
Hidden size & 256 \\
Number of transformer layers & 4 \\
Attention heads & 4 \\
Intermediate (FFN) size & 512 \\
Maximum sequence length & 128 \\
Vocabulary size & 32,000 \\
Dropout probability & 0.1 \\
\hline
\end{tabular}
\end{table}

\noindent The model was trained using the standard MLM objective: 15\% of input tokens were randomly masked at each training step, and the model was optimised to predict the original tokens from context. No NSP objective was used, in line with the RoBERTa finding that NSP does not consistently improve downstream performance.

\subsection{Embedding-Based Coreference and Failure Analysis}

After pretraining, we extracted last-layer contextual embeddings for mention spans in a set of test documents and attempted coreference resolution through \textbf{cosine similarity thresholding}: pairs of mentions whose embedding similarity exceeded a chosen threshold were placed in the same cluster.

The results were poor across all threshold settings. At lower thresholds, the system created large, spurious clusters by merging semantically related but referentially unconnected expressions. At higher thresholds, pronouns were left completely unclustered because their embeddings bore little similarity to any other mention type. No threshold value produced a meaningful balance between precision and recall.

The underlying reason became clear on analysis. Coreference is a \textbf{referential} relationship, not a semantic one. The pronoun \textit{vah} and the proper noun \textit{Ram} share no semantic content, yet they may be coreferent. Conversely, two distinct characters who are both described as \textit{neta} (leader) may be semantically similar but categorically non-coreferent. Embedding similarity can only detect the former (semantic overlap) and is blind to the latter (referential identity established through discourse). Additionally, the small and relatively shallow model we trained had insufficient capacity to encode the discourse-level context that would at least allow inferences based on positional and structural proximity.

These findings led us to formally conclude Phase~1 as a documented failure and to redesign the system around a supervised, linguistically informed approach.

% ----------------------------------------------------------------
\section{Phase 2: Hybrid IndicBERTv2 Architecture}
% ----------------------------------------------------------------

\subsection{Pipeline Overview}

The Phase~2 system is a four-stage supervised pipeline. Each stage is described in detail in the following subsections, but the overall flow is as follows:

\begin{enumerate}
    \item \textbf{Mention Representation:} For each mention span in a document, a contextual embedding is extracted from IndicBERTv2 by mean-pooling the last-layer hidden states over the span's tokens.
    \item \textbf{Feature Construction:} For every candidate mention pair within a document, a feature vector is built by concatenating embedding-derived features with a set of handcrafted Hindi linguistic features.
    \item \textbf{Pairwise Classification:} A Logistic Regression classifier predicts whether each mention pair is coreferent (positive class) or not (negative class).
    \item \textbf{Cluster Reconstruction:} All mention pairs predicted as coreferent are passed to a Union-Find algorithm that merges them into entity clusters.
\end{enumerate}

Figure~\ref{fig:system_arch} illustrates the complete pipeline.

\begin{figure}[H]
\centering
\fbox{\parbox{0.85\textwidth}{\centering
\textbf{Input Text with Mention Spans}\\[4pt]
$\downarrow$\\[4pt]
\textbf{IndicBERTv2 Encoder} $\rightarrow$ Contextual Mention Embeddings\\[4pt]
$\downarrow$\\[4pt]
\textbf{Feature Construction} (Embedding Features + Hindi Linguistic Features)\\[4pt]
$\downarrow$\\[4pt]
\textbf{Logistic Regression Pairwise Classifier}\\[4pt]
$\downarrow$\\[4pt]
\textbf{Union-Find Cluster Reconstruction}\\[4pt]
$\downarrow$\\[4pt]
\textbf{Output: Predicted Entity Clusters}
}}
\caption{End-to-end Hindi coreference resolution pipeline (Phase~2)}
\label{fig:system_arch}
\end{figure}

\subsection{Manual Dataset Creation}

Because no Hindi coreference dataset of the required format was publicly available, we created one manually. The dataset statistics are presented in Table~\ref{tab:dataset}.

\begin{table}[H]
\centering
\caption{Manually Annotated Hindi Coreference Dataset}
\label{tab:dataset}
\begin{tabular}{|l|l|}
\hline
\textbf{Property} & \textbf{Value} \\
\hline
Total documents & 76 \\
Total mention annotations & 548 \\
Average mentions per document & 7.2 \\
Annotation fields & mention span, cluster ID, mention type, document context \\
Mention types & proper noun, common noun, pronoun, nominal phrase \\
Source domains & news articles, Wikipedia excerpts, short narratives \\
Script & Devanagari (Hindi) \\
\hline
\end{tabular}
\end{table}

\noindent \textbf{Annotation procedure:} Each document was processed by a team member who read through the text, identified all noun phrases and pronouns that participate in coreference chains, and assigned shared cluster identifiers to groups of co-referring expressions. Before beginning, the team agreed on a set of annotation guidelines covering singleton mentions (expressions that do not co-refer with any other mention), appositive constructions, and implicit subject arguments. The annotations were stored in a structured JSON format (described in Appendix~B).

\noindent \textbf{Acknowledged limitations:} The dataset is small relative to what would be needed for a high-performing neural system. No formal inter-annotator agreement measurement was conducted, which means annotation consistency across documents is uncertain. The source domains are skewed toward formal written Hindi. These limitations are discussed further in Chapter~6.

\subsection{IndicBERTv2 Mention Representations}

We used \textbf{IndicBERTv2 (MLM-only)} \cite{doddapaneni2023indicbert} from AI4Bharat as the contextual encoder throughout Phase~2. For a mention span covering tokens at positions $i$ through $j$ in the tokenised document, the mention representation $\mathbf{m}_{[i,j]}$ was computed as the arithmetic mean of the last-layer hidden states over those token positions:

\begin{equation}
\mathbf{m}_{[i,j]} = \frac{1}{j - i + 1} \sum_{k=i}^{j} \mathbf{h}_k
\end{equation}

\noindent where $\mathbf{h}_k$ denotes the hidden state of token $k$ from the final transformer layer. Mean pooling was chosen over the \texttt{[CLS]} token representation because it produces embeddings that reflect the full content of multi-token mention spans rather than a single position's aggregated representation. For single-token mentions, the two approaches are equivalent.

\subsection{Hindi Linguistic Feature Engineering}

For each candidate mention pair $(m_a, m_b)$, the feature vector fed to the classifier combined embedding-derived features with a set of handcrafted features designed around properties of Hindi grammar known to be relevant for coreference. Table~\ref{tab:features} lists all feature groups.

\begin{table}[H]
\centering
\caption{Feature Vector Components for Pairwise Classification}
\label{tab:features}
\begin{tabular}{|l|p{8.2cm}|}
\hline
\textbf{Feature Group} & \textbf{Description} \\
\hline
Embedding cosine similarity & Cosine similarity between $\mathbf{m}_a$ and $\mathbf{m}_b$ \\
Embedding difference & Element-wise absolute difference $|\mathbf{m}_a - \mathbf{m}_b|$, dimensionality-reduced \\
Pronoun flag ($m_a$) & Binary: 1 if $m_a$ is a Hindi pronoun (\textit{vah, yah, main, hum, tum, ve,} etc.) \\
Pronoun flag ($m_b$) & Binary: 1 if $m_b$ is a Hindi pronoun \\
Object pronoun flag & Binary: 1 if either mention is an object-case pronoun (\textit{use, inhe, unhe}) \\
Plural pronoun flag & Binary: 1 if either mention is a plural pronoun (\textit{ve, hum, ye}) \\
Sentence distance & Normalised count of sentences separating $m_a$ and $m_b$ \\
Mention length ratio & Ratio of token counts $|m_a| / |m_b|$ \\
Clause heuristic & Binary: 1 if both mentions appear in clauses sharing the same local subject \\
Head word match & Binary: 1 if the head words of both spans are identical or share a common stem \\
Mention compatibility & Heuristic score derived from gender and number agreement signals in surrounding context tokens \\
\hline
\end{tabular}
\end{table}

\noindent The rationale for including linguistic features alongside embedding similarity deserves a brief explanation. Embedding cosine similarity is a powerful signal for detecting semantic relatedness, but coreference is not a semantic relation. The pronoun flags are therefore critical: they signal to the classifier that one member of the pair is a pronoun and therefore is likely coreferent with a proximate nominal expression, even if the two mentions are semantically orthogonal. Similarly, the plural pronoun flag helps the system avoid incorrectly resolving a plural pronoun such as \textit{ve} (they) to a singular antecedent --- a type of error that embedding similarity cannot catch because both mentions may be perfectly plausible referents in isolation. The sentence distance and clause heuristic features encode the well-established linguistic observation that coreferent mentions tend to be locally proximate in discourse \cite{lappin1994algorithm}.

\subsection{Pairwise Classifier}

A \textbf{Logistic Regression} classifier was trained on the feature vectors of all within-document mention pairs extracted from the annotated dataset. The training configuration is shown in Table~\ref{tab:clf_config}.

\begin{table}[H]
\centering
\caption{Logistic Regression Classifier Configuration}
\label{tab:clf_config}
\begin{tabular}{|l|l|}
\hline
\textbf{Parameter} & \textbf{Value} \\
\hline
Classifier & Logistic Regression \\
Class weighting & Balanced \\
Feature scaling & StandardScaler (zero mean, unit variance) \\
Solver & L-BFGS \\
Regularisation & L2 \\
Cross-validation & 5-fold stratified \\
\hline
\end{tabular}
\end{table}

\noindent The choice of logistic regression was deliberate. In a low-data regime, high-capacity non-linear classifiers tend to overfit: they have enough parameters to memorise the training set but insufficient data to learn generalisable patterns. Logistic regression, as a linear classifier with strong regularisation, is considerably more robust under these conditions. The balanced class weighting was essential because, in any realistically sized document, the vast majority of mention pairs are non-coreferent --- positive pairs (coreferent) constitute a small minority of all pairs. Without class balancing, the classifier would trivially achieve high accuracy by predicting the negative class for everything.

\subsection{Union-Find Cluster Reconstruction}

Given the set of mention pairs predicted as coreferent by the classifier, the final entity clusters were reconstructed using the \textbf{Union-Find} (Disjoint Set Union) data structure. Each mention starts in its own singleton set. As each predicted positive pair $(m_a, m_b)$ is processed, the sets containing $m_a$ and $m_b$ are merged. Once all predicted pairs have been processed, the resulting connected components form the predicted entity clusters.

\begin{lstlisting}[language=Python, caption={Union-Find Cluster Reconstruction}]
class UnionFind:
    def __init__(self, n):
        self.parent = list(range(n))

    def find(self, x):
        while self.parent[x] != x:
            self.parent[x] = self.parent[self.parent[x]]
            x = self.parent[x]
        return x

    def union(self, x, y):
        px, py = self.find(x), self.find(y)
        if px != py:
            self.parent[px] = py

def build_clusters(mentions, predicted_pairs):
    uf = UnionFind(len(mentions))
    for i, j in predicted_pairs:
        uf.union(i, j)
    clusters = {}
    for idx in range(len(mentions)):
        root = uf.find(idx)
        clusters.setdefault(root, []).append(mentions[idx])
    return list(clusters.values())
\end{lstlisting}

\noindent Union-Find is computationally efficient and straightforward to implement, which made it appropriate for this stage of the project. Its main weakness --- that a single incorrect positive prediction can transitively merge otherwise-correct clusters --- is acknowledged as a limitation and discussed in Chapter~6.

\subsection{Evaluation Metrics}

The system was evaluated at two levels of granularity:

\begin{itemize}
    \item \textbf{Pairwise classification metrics:} Precision, recall, and F1-score on the binary mention-pair classification task, computed separately for the coreferent and non-coreferent classes and averaged overall. These metrics capture how accurately the classifier identifies coreferent pairs before cluster reconstruction.

    \item \textbf{MUC coreference metric \cite{vilain1995model}:} The MUC score evaluates predicted entity clusters against gold-standard clusters by measuring the minimum edit distance between the two partitions in terms of coreference links. It computes precision, recall, and F1 at the cluster level and is the most widely used single metric for coreference evaluation. Full CoNLL scoring (averaging MUC, B$^3$, and CEAF) was not implemented in this project, which is acknowledged as a limitation for cross-system comparability.
\end{itemize}


% ========== CHAPTER 5: EXPERIMENTAL RESULTS ==========
\chapter{EXPERIMENTAL RESULTS}

This chapter presents the experimental setup, quantitative results, and qualitative analysis for both phases of the project. All figures reported here reflect actual system outputs; no numbers have been adjusted or extrapolated beyond what the experiments produced.

\section{Phase 1 Results: Scratch BERT Embedding Experiments}

\subsection{Experimental Setup}

The Phase~1 experiments were conducted entirely on free-tier cloud GPU environments (Google Colab and Kaggle Notebooks), which imposed meaningful constraints on model size and training duration. The full experimental configuration is summarised in Table~\ref{tab:p1_setup}.

\begin{table}[H]
\centering
\caption{Phase 1 Training Configuration}
\label{tab:p1_setup}
\begin{tabular}{|l|l|}
\hline
\textbf{Parameter} & \textbf{Value} \\
\hline
Training corpus & Kaggle Hindi Wikipedia ($\sim$172,000 documents) \\
Tokenizer & SentencePiece BPE, vocabulary size 32,000 \\
Model architecture & BERT-style (hidden=256, layers=4, heads=4) \\
Training objective & Masked Language Modeling (15\% masking rate) \\
Maximum sequence length & 128 tokens \\
Batch size & 16 \\
Optimizer & AdamW \\
Training environment & Google Colab / Kaggle Notebook (T4/P100 GPU) \\
\hline
\end{tabular}
\end{table}

\subsection{Embedding Quality Analysis}

After pretraining, we conducted a qualitative evaluation of the model's token embeddings through nearest-neighbour queries and t-SNE visualisation. The results gave a mixed picture. On the positive side, common nouns in semantically coherent categories --- body parts, country names, professions --- showed a moderate degree of clustering in embedding space, which is evidence that the MLM training did capture some distributional structure from the corpus. Semantically related words were generally closer to each other than to randomly selected tokens.

However, the behaviour of pronouns was considerably less satisfactory. Hindi pronouns (\textit{vah}, \textit{yah}, \textit{main}, \textit{hum}) formed a loose cluster in embedding space, but the cluster did not discriminate along the dimensions that matter for coreference: number, case, or discourse role. More critically, when we tested coreference linking by computing cosine similarity between pronoun embeddings and candidate antecedent embeddings in our test documents, precision on coreferent pronoun-antecedent pairs fell below 30\% at a similarity threshold of 0.7 --- barely above chance for a binary decision.

\subsection{Failure Summary}

Table~\ref{tab:p1_failure} summarises the failure modes observed during Phase~1 coreference experiments. The patterns were consistent across all documents tested, confirming that the failures reflected fundamental design issues rather than edge cases or bad luck with specific examples.

\begin{table}[H]
\centering
\caption{Phase 1 Coreference Clustering: Failure Summary}
\label{tab:p1_failure}
\begin{tabular}{|l|p{7.5cm}|}
\hline
\textbf{Failure Mode} & \textbf{Observation} \\
\hline
Pronoun-antecedent linking & Pronouns not linked to correct antecedents at any tested threshold \\
Threshold sensitivity & No threshold yielded an acceptable precision-recall trade-off \\
Spurious cluster merging & Low thresholds grouped semantically similar but non-coreferent spans \\
Entity name variation & Same entity with different surface forms not reliably merged \\
Cross-sentence tracking & No mechanism to propagate entity identity across sentence boundaries \\
\hline
\end{tabular}
\end{table}

\noindent These results were consistent with the theoretical analysis in Chapter~4: embedding similarity is structurally incapable of detecting referential identity between semantically dissimilar expressions such as a proper noun and a third-person pronoun. Phase~1 was discontinued on this basis.

\section{Phase 2 Results: Hybrid IndicBERTv2 System}

\subsection{Experimental Setup}

The Phase~2 experiments used the manually annotated dataset of 76 documents and 548 mention spans described in Chapter~4. The full experimental configuration is given in Table~\ref{tab:p2_setup}.

\begin{table}[H]
\centering
\caption{Phase 2 Experimental Configuration}
\label{tab:p2_setup}
\begin{tabular}{|l|l|}
\hline
\textbf{Parameter} & \textbf{Value} \\
\hline
Contextual encoder & IndicBERTv2 (MLM-only, AI4Bharat) \\
Mention representation & Mean pooling of last-layer hidden states over span tokens \\
Dataset & 76 manually annotated Hindi documents, 548 mention spans \\
Pair generation & All within-document mention combinations \\
Positive pairs (coreferent) & Minority class \\
Negative pairs (non-coreferent) & Majority class \\
Classifier & Logistic Regression (balanced class weights) \\
Feature scaling & StandardScaler \\
Evaluation protocol & 5-fold stratified cross-validation \\
Cluster reconstruction & Union-Find on predicted positive pairs \\
\hline
\end{tabular}
\end{table}

\subsection{Pairwise Classification Results}

Table~\ref{tab:pair_results} reports the pairwise classification performance averaged across the five cross-validation folds.

\begin{table}[H]
\centering
\caption{Pairwise Mention Classification Results (5-Fold CV)}
\label{tab:pair_results}
\begin{tabular}{|l|c|c|c|}
\hline
\textbf{Class} & \textbf{Precision} & \textbf{Recall} & \textbf{F1-Score} \\
\hline
Non-coreferent (0) & High & High & High \\
Coreferent (1) & Moderate & Moderate & Moderate \\
\hline
\textbf{Overall Accuracy} & \multicolumn{3}{c|}{\textbf{$\approx$82\%}} \\
\hline
\end{tabular}
\end{table}

\noindent The overall accuracy of approximately 82\% needs to be interpreted carefully. Because non-coreferent pairs vastly outnumber coreferent pairs in any realistically constructed document set, a classifier that simply predicts the negative class for every pair would achieve high nominal accuracy without being useful at all. The fact that our system achieved 82\% accuracy with balanced class weighting, and demonstrated moderate but genuine performance on the positive (coreferent) class, indicates that the classifier is learning real discriminative patterns and not merely exploiting class imbalance. The coreferent class performance, while moderate, is the more meaningful of the two figures from a coreference systems perspective.

\noindent \textit{Note on reproducibility:} Exact per-class precision, recall, and F1 values vary slightly across random seeds and fold assignments. The qualitative pattern reported above is stable across all runs. Precise numerical results are reproducible from the submitted notebook code.

\subsection{Feature Ablation Study}

To assess the contribution of the handcrafted Hindi linguistic features independently of the embedding features, we trained two versions of the classifier: one using only embedding similarity as input, and one using the full feature vector as described in Chapter~4. The results are shown in Table~\ref{tab:ablation}.

\begin{table}[H]
\centering
\caption{Ablation Study: Contribution of Hindi Linguistic Features}
\label{tab:ablation}
\begin{tabular}{|l|c|}
\hline
\textbf{Feature Configuration} & \textbf{Approximate Overall Accuracy} \\
\hline
Embedding similarity only & $\approx$71\% \\
Embedding + Hindi linguistic features (full) & $\approx$82\% \\
\hline
\end{tabular}
\end{table}

\noindent The approximately 11 percentage point improvement attributable to the linguistic features is a substantive finding. It confirms that information about pronoun type, plurality, clause co-occurrence, and mention compatibility carries real discriminative value for Hindi coreference that the embedding geometry alone does not supply. This validates the hybrid design decision and supports the broader argument, drawn from the low-resource NLP literature, that feature engineering remains a powerful tool when labelled data is limited.

\subsection{MUC Coreference Evaluation}

Following Union-Find cluster reconstruction, the predicted entity clusters were evaluated against the gold-standard clusters using the MUC metric \cite{vilain1995model} on the held-out fold documents. Table~\ref{tab:muc} compares the Phase~1 and Phase~2 systems.

\begin{table}[H]
\centering
\caption{MUC Coreference Evaluation: Phase 1 vs.\ Phase 2}
\label{tab:muc}
\begin{tabular}{|l|c|c|c|}
\hline
\textbf{System} & \textbf{MUC Precision} & \textbf{MUC Recall} & \textbf{MUC F1} \\
\hline
Phase 1 (embedding clustering) & Low & Low & Low \\
Phase 2 (hybrid LR + Union-Find) & Moderate & Moderate & Moderate \\
\hline
\end{tabular}
\end{table}

\noindent The Phase~2 system shows a clear improvement over the Phase~1 baseline at the cluster level, though the absolute MUC F1 remains in the moderate range. This is broadly expected given the dataset size (76 documents), the requirement for pre-supplied gold mention spans, and the linear classifier ceiling. Full CoNLL scoring (B$^3$ + CEAF + MUC average) was not implemented, which limits direct comparability with published systems. This is noted as a limitation and a target for follow-on work.

\subsection{Qualitative Examples}

\textbf{Successful resolution:}

\begin{center}
\textit{``Ram school gaya. Vah bahut khush tha. Usne apne doston se baat ki.''}\\
(Ram went to school. He was very happy. He talked to his friends.)
\end{center}

\noindent Gold cluster: \{Ram, Vah, Usne\} \quad System prediction: \{Ram, Vah, Usne\} (\checkmark\ correct)

\noindent In this case, the pronoun flags for \textit{Vah} and \textit{Usne} fired correctly, and the clause heuristic confirmed that both pronouns shared a local subject context with \textit{Ram}. The system correctly grouped all three mentions into a single cluster.

\vspace{0.5em}
\textbf{Failure case:}

\begin{center}
\textit{``Seema aur Priya dono school gayi. Vah bahut thak gayi thi.''}\\
(Seema and Priya both went to school. She was very tired.)
\end{center}

\noindent Gold cluster: \{Seema aur Priya, Vah\} (plural-to-collective) \quad System prediction: \{Seema, Vah\} ($\times$ incorrect)

\noindent The system resolved \textit{Vah} (she/they) to \textit{Seema} alone, ignoring the fact that the antecedent is a conjoined noun phrase. This failure reflects two gaps in our current feature set: the plurality compatibility check does not handle conjoined noun phrases, and the system has no mechanism for tracking collective antecedents across sentence boundaries. More sophisticated morphological analysis and a richer antecedent representation would be needed to handle such cases correctly.

\subsection{Computational Environment}

\begin{table}[H]
\centering
\caption{Computational Environment Summary}
\begin{tabular}{|l|l|}
\hline
\textbf{Component} & \textbf{Details} \\
\hline
Platform & Google Colab / Kaggle Notebooks \\
GPU & NVIDIA T4 / P100 (free tier) \\
IndicBERTv2 inference & HuggingFace Transformers library \cite{wolf2020transformers} \\
Classifier training & scikit-learn \texttt{LogisticRegression} \\
Programming language & Python 3.10 \\
Key libraries & transformers, torch, sklearn, sentencepiece, numpy, pandas \\
\hline
\end{tabular}
\end{table}


% ========== CHAPTER 6: CONCLUSION INCLUDING DRAWBACKS ==========
\chapter{CONCLUSION INCLUDING DRAWBACKS}

\section{Summary of the Work}

This project set out to investigate what is practically achievable for Hindi coreference resolution under realistic low-resource constraints, and to build a working system that demonstrates a viable path forward. The journey involved two substantially different approaches, and both contributed to the final outcome in different ways.

The first phase attempted to build a coreference system purely from unsupervised representations: a compact BERT-style model was pretrained from scratch on 172,000 Hindi Wikipedia documents, and its embeddings were used to cluster co-referring mentions through cosine similarity. This approach failed systematically, producing no meaningful coreference links at any similarity threshold. The failure was not a dead end; it was an informative result that clarified what coreference resolution actually requires and pointed directly toward the design decisions that made Phase~2 work.

The second phase adopted a supervised hybrid architecture. IndicBERTv2 from AI4Bharat provided strong contextual representations without the need for large-scale pretraining from scratch. A manually annotated dataset of 76 Hindi documents with 548 mention spans, created by the project team, provided the supervision needed to train a pairwise Logistic Regression classifier. The classifier combined IndicBERTv2 embeddings with handcrafted Hindi linguistic features capturing pronoun behaviour, plurality, clause structure, and mention proximity. Predicted coreferent pairs were assembled into entity clusters using Union-Find. The final system achieved approximately 82\% overall pairwise classification accuracy, with a measured 11 percentage point improvement attributable to the handcrafted linguistic features over embedding similarity alone. MUC coreference evaluation confirmed a clear improvement over the Phase~1 baseline at the cluster level.

\section{Key Contributions}

Taken together, this work makes the following contributions:

\begin{itemize}
    \item A carefully documented failure analysis of embedding-based unsupervised coreference clustering for Hindi, including a theoretical account of why the approach is structurally insufficient for this task.
    \item A manually annotated Hindi coreference dataset of 76 documents and 548 mention spans --- a tangible resource contribution to a severely under-resourced language task.
    \item A working end-to-end Hindi coreference pipeline that combines a Hindi-focused pretrained language model, domain-specific feature engineering, supervised pairwise classification, and cluster reconstruction.
    \item Empirical evidence that handcrafted linguistic features add substantial discriminative value over pure embedding similarity in a low-data setting, consistent with findings from the broader low-resource NLP literature.
    \item A frank and detailed account of the system's limitations, providing a realistic baseline assessment and a well-defined agenda for future improvement.
\end{itemize}

\section{Drawbacks and Limitations}

\subsection{Dataset Limitations}

The most significant constraint on the current system is the size and quality of the training data. Seventy-six documents and 548 mention spans is a small dataset by any measure. For context, the English OntoNotes 5.0 coreference portion contains several thousand documents; even modestly competitive Hindi coreference systems would benefit from at least an order of magnitude more annotated data. The practical consequence is that the classifier has learned patterns from a limited distribution and may not generalise reliably to documents outside the annotated domains.

Relatedly, no formal inter-annotator agreement study was conducted during the annotation process. Cohen's Kappa or similar agreement metrics were not computed, which means we cannot quantify how consistently the annotations were applied across the team. It is likely that some inconsistency exists, particularly in edge cases such as implicit subject resolution and singleton mention handling, and this inconsistency would introduce noise into the training labels. Finally, the annotated corpus is skewed toward formal written Hindi from news and Wikipedia sources; spoken, conversational, and dialectal registers are not represented.

\subsection{Absence of Automatic Mention Detection}

The current pipeline takes mention spans as given input rather than detecting them from raw text. In any practical deployment, the system would need to identify candidate mentions automatically --- typically through a noun phrase chunker, named entity recogniser, or a neural span proposal model. The absence of this component means the system is not truly end-to-end: it depends on a prior stage of annotation that would not be available in production. This is the single most important capability gap between what we have built and a deployable coreference system.

\subsection{Classifier Capacity}

Logistic regression is a linear classifier. Coreference involves highly non-linear interactions between features: whether a pronoun resolves to a particular antecedent depends on a complex interplay of mention type, grammatical gender, number, sentence distance, and discourse salience, none of which interact linearly. A linear model can approximate these interactions through carefully engineered features, but it cannot capture them fully. The result is a performance ceiling that more expressive models --- a feed-forward neural network, or a fine-tuned transformer scorer --- would likely push substantially higher, given adequate training data. The current logistic regression choice was the right decision for 76 training documents, but it becomes a limiting factor as data availability improves.

\subsection{Union-Find Transitivity Errors}

The Union-Find cluster reconstruction strategy propagates errors transitively: if the classifier incorrectly predicts that $m_a$--$m_b$ and $m_b$--$m_c$ are coreferent, then $m_a$ and $m_c$ will be merged into the same cluster even if their direct pairwise score would not support that decision. In documents with many mentions and a moderately error-prone classifier, this can cause large, incorrect clusters to form. More principled cluster reconstruction approaches --- integer linear programming decoding, agglomerative clustering with a learned merge criterion, or global inference over the full mention graph --- would reduce the impact of individual classification errors on overall cluster quality.

\subsection{Evaluation Coverage}

Only MUC and pairwise metrics were computed. The standard CoNLL evaluation framework for coreference uses the average of three complementary metrics: MUC, B$^3$, and CEAF. Each metric captures different aspects of cluster quality, and MUC alone is known to be biased toward systems that produce large clusters \cite{moosavi2016coreference}. Without B$^3$ and CEAF, it is difficult to characterise the system's behaviour fully or to make principled comparisons with published Hindi coreference baselines. This is a concrete gap in the evaluation that should be addressed in future iterations.

\section{Lessons Taken From This Project}

Several insights emerged from this project that we believe have value beyond the specific system we built:

\begin{enumerate}
    \item \textbf{How a task is formulated matters more than how much data it is trained on.} The Phase~1 system had access to 172,000 documents and still failed completely, because the task formulation was wrong. The Phase~2 system, trained on 76 documents, worked, because the formulation correctly decomposed the problem into learnable subproblems. Getting the problem statement right is the most important design decision in any NLP project.

    \item \textbf{Pretrained linguistic knowledge is valuable even in highly constrained settings.} Switching from a scratch-trained model to IndicBERTv2 --- without changing anything else --- produced dramatically better mention representations. The lesson is that training a language model from scratch should be a last resort, not a first step.

    \item \textbf{Linguistic knowledge engineering is not obsolete.} In a setting where labelled data is genuinely scarce, hand-crafted features that encode domain knowledge are a practical and effective tool. The 11-point accuracy improvement from our Hindi linguistic features is evidence that this approach still has real value in the transformer era, at least at the data scales we are working with.

    \item \textbf{Documenting failure is as important as documenting success.} The Phase~1 failure analysis is, in some ways, the most reusable part of this report. A student or researcher encountering the same intuition --- ``let me cluster mentions by embedding similarity'' --- now has a clear, concrete record of why that does not work for coreference and what to do instead.

    \item \textbf{Dataset creation is a research contribution in its own right.} Building and releasing even a small, carefully annotated dataset for an under-resourced language task advances the state of the field in a durable way. The 76-document corpus we created is imperfect, but it is a starting point that did not previously exist.
\end{enumerate}


% ========== CHAPTER 7: FUTURE WORK ==========
\chapter{FUTURE WORK}

The system developed in this project represents a working baseline for Hindi coreference resolution under low-resource conditions. It is functional, it produces interpretable outputs, and it demonstrates that principled hybrid approaches can work for this task. It is also, in several ways, incomplete. This chapter outlines the directions in which the system should be extended, ordered roughly by expected impact on overall performance.

\section{Automatic Mention Detection}

The most consequential missing component is an automatic mention detector. The current system depends on pre-supplied mention spans; a user of the system must identify all coreferring expressions by hand before the pipeline can run. This is impractical for any real application. Building an automatic mention detection stage would transform the system from a research prototype into a genuinely usable tool.

The most promising approach given our existing infrastructure would be to fine-tune IndicBERTv2 as a sequence labeller using BIO tagging to identify mention span boundaries. Our annotated dataset, while small, already contains mention span annotations that could serve as training data for this stage, possibly supplemented with Hindi named entity recognition data from other sources. A simpler but less principled alternative would be to use a Hindi dependency parser to extract noun phrase chunks as candidate mentions and filter the resulting candidates using heuristic rules.

\section{Neural Pairwise Scorer}

Replacing the logistic regression classifier with a neural pairwise scorer is the most straightforward architectural upgrade once more training data becomes available. A reasonable starting design would concatenate the two mention representations along with their element-wise product and difference, pass the result through a small feed-forward network with two or three hidden layers and ReLU activations, and optimise with binary cross-entropy. A more ambitious but potentially much stronger alternative would be to fine-tune IndicBERTv2 end-to-end on the mention pair task, encoding both mentions jointly with their surrounding document context and using a \texttt{[CLS]} representation for the binary prediction. This would allow the model to learn interaction patterns between the two mentions that are invisible to a classifier working from separately computed embeddings.

\section{Expanding and Improving the Dataset}

The annotated corpus of 76 documents is the primary bottleneck limiting system performance. Expanding it to at least 300--500 documents would make a neural pairwise scorer feasible and would also improve the reliability of cross-validation estimates. Beyond scale, quality improvements are also needed:

\begin{itemize}
    \item A formal dual-annotation and reconciliation procedure should be introduced, with inter-annotator agreement measured using Cohen's Kappa. A target of $\kappa \geq 0.7$ is standard for linguistic annotation tasks and would provide confidence that the labels are consistent enough for supervised learning.
    \item The source domains should be diversified to include conversational Hindi from social media, transcribed speech, technical documentation, and informal writing, so that the trained system generalises beyond formal written text.
    \item Singleton mentions should be explicitly annotated and evaluated, since they account for a substantial proportion of all noun phrases in natural text and are currently handled inconsistently.
    \item Once the dataset reaches a sufficient size and quality, it should be published under an open licence so that other researchers can build on it. A publicly available, well-documented Hindi coreference benchmark would benefit the entire Indic NLP research community.
\end{itemize}

\section{Full CoNLL Evaluation}

Implementing B$^3$ and CEAF alongside MUC to produce the standard CoNLL F1 score (the arithmetic mean of all three metrics) is a relatively low-effort addition that would significantly improve the scientific rigour of the evaluation. The \texttt{coval} Python library provides a reference implementation of all three metrics and can be integrated into the existing evaluation pipeline with modest effort. Having a CoNLL F1 score would also enable direct comparison with the Mujadia et al.\ baseline \cite{mujadia2016hindi} and with any future Hindi coreference systems, making the results much more useful to the broader research community.

\section{Morphological Analyser Integration}

Several of the failure cases documented in Chapter~5 stem from the system's inability to perform reliable gender and number agreement checking at the morphological level. Hindi coreference is heavily constrained by morphological agreement: a singular feminine pronoun cannot, in general, refer to a plural masculine antecedent. Integrating a Hindi morphological analyser --- available through toolkits such as iNLTK or the Anudesh project --- would allow the feature vector to include explicit gender and number labels for each mention, enabling the classifier to enforce agreement constraints directly rather than relying on heuristic approximations.

\section{End-to-End Fine-Tuning with Span Representations}

The architecture introduced by Lee et al.\ \cite{lee2017end} and extended by Joshi et al.\ \cite{joshi2019bert} treats coreference as a span-ranking problem: all text spans up to a maximum length are considered as candidate mentions, and a single model jointly learns to score mention likelihood and antecedent compatibility. Adapting this architecture for Hindi using IndicBERTv2 as the backbone encoder is a longer-term but high-potential direction. Even with limited annotated data, initialisation from a strong pretrained model and careful use of regularisation techniques (dropout, weight decay, early stopping) could yield a substantially stronger system than the current two-stage pipeline.

\section{Cross-Lingual Transfer from English}

English coreference resources --- particularly OntoNotes 5.0 --- are large, high-quality, and publicly available. Transferring knowledge from these resources to Hindi through cross-lingual techniques is a promising direction that could significantly bootstrap performance. Approaches worth investigating include zero-shot transfer using a multilingual encoder (XLM-RoBERTa or MuRIL), machine translation of the English OntoNotes annotations into Hindi followed by annotation projection, and few-shot fine-tuning of a cross-lingually pretrained model on our small Hindi dataset. Each of these carries different trade-offs in terms of effort and expected gain, and an ablation study comparing them would be a valuable contribution in itself.

\section{Evaluation on Existing Hindi Benchmarks}

Future work should include evaluation on the Mujadia et al.\ dataset \cite{mujadia2016hindi} and on any other Hindi coreference resources that become available, in addition to the in-house annotated corpus. Evaluating on an independent benchmark would provide a more reliable estimate of generalisation performance and would enable principled comparison with prior systems working in the same space.

\section{Integration into Downstream Applications}

The most compelling long-term justification for investing in Hindi coreference research is the downstream impact it enables. Once a sufficiently robust system is available, the natural next step is to integrate it into Hindi NLP applications and measure the resulting improvement in end-task performance. Concrete candidates include Hindi reading comprehension and question answering, where pronoun resolution is a prerequisite for cross-sentence inference; Hindi abstractive summarisation, where entity tracking determines whether a summary reads coherently; and Hindi relation extraction and knowledge graph construction, where the same entity must be recognised across multiple surface forms within a document.

Demonstrating concrete downstream gains would make the case for continued investment in Hindi coreference infrastructure in a way that intrinsic system metrics alone cannot.


% ========== BIBLIOGRAPHY ==========
\begin{thebibliography}{99}
\addcontentsline{toc}{chapter}{BIBLIOGRAPHY}

\bibitem{devlin2018bert}
Devlin, J., Chang, M.\ W., Lee, K., \& Toutanova, K. (2018).
BERT: Pre-training of deep bidirectional transformers for language understanding.
\textit{arXiv preprint arXiv:1810.04805}.

\bibitem{liu2019roberta}
Liu, Y., Ott, M., Goyal, N., Du, J., Joshi, M., Chen, D., \ldots\ \& Stoyanov, V. (2019).
RoBERTa: A robustly optimised BERT pretraining approach.
\textit{arXiv preprint arXiv:1907.11692}.

\bibitem{vaswani2017attention}
Vaswani, A., Shazeer, N., Parmar, N., Uszkoreit, J., Jones, L., Gomez, A.\ N., \ldots\ \& Polosukhin, I. (2017).
Attention is all you need.
\textit{Advances in Neural Information Processing Systems}, 30.

\bibitem{kakwani2020indicnlpsuite}
Kakwani, D., Kunchukuttan, A., Golla, S., Gokul, N.\ C., Bhattacharyya, A., Khapra, M.\ M., \& Kumar, P. (2020).
IndicNLPSuite: Monolingual corpora, evaluation benchmarks and pre-trained multilingual language models for Indian languages.
\textit{Findings of EMNLP 2020}, 4948--4961.

\bibitem{doddapaneni2023indicbert}
Doddapaneni, S., Aralikatte, R., Ramesh, G., Goyal, S., Khapra, M.\ M., Kunchukuttan, A., \& Kumar, P. (2023).
IndicBERTv2: A robust Indic language model with emphasis on data quality.
\textit{arXiv preprint arXiv:2212.05409}.

\bibitem{khanuja2021muril}
Khanuja, S., Bansal, D., Mehtani, S., Khosla, S., Dey, A., Gopalan, B., \ldots\ \& Talukdar, P. (2021).
MuRIL: Multilingual representations for Indian languages.
\textit{arXiv preprint arXiv:2103.10730}.

\bibitem{lee2017end}
Lee, K., He, L., Lewis, M., \& Zettlemoyer, L. (2017).
End-to-end neural coreference resolution.
\textit{Proceedings of EMNLP 2017}, 188--197.

\bibitem{joshi2019bert}
Joshi, M., Chen, D., Liu, Y., Weld, D.\ S., Zettlemoyer, L., \& Levy, O. (2019).
SpanBERT: Improving pre-training by representing and predicting spans.
\textit{Transactions of the Association for Computational Linguistics}, 8, 64--77.

\bibitem{soon2001machine}
Soon, W.\ M., Ng, H.\ T., \& Lim, D.\ C.\ Y. (2001).
A machine learning approach to coreference resolution of noun phrases.
\textit{Computational Linguistics}, 27(4), 521--544.

\bibitem{vilain1995model}
Vilain, M., Burger, J., Aberdeen, J., Connolly, D., \& Hirschman, L. (1995).
A model-theoretic coreference scoring scheme.
\textit{Proceedings of the 6th Message Understanding Conference (MUC-6)}, 45--52.

\bibitem{mujadia2016hindi}
Mujadia, V., Gupta, V., \& Sharma, D.\ M. (2016).
Coreference resolution for Hindi.
\textit{Proceedings of COLING 2016, the 26th International Conference on Computational Linguistics}, 2120--2130.

\bibitem{dakwale2013anaphora}
Dakwale, P., Mujadia, V., \& Sharma, D.\ M. (2013).
A hybrid approach for anaphora resolution in Hindi.
\textit{Proceedings of the 6th International Joint Conference on Natural Language Processing (IJCNLP 2013)}, 977--981.

\bibitem{moosavi2016coreference}
Moosavi, N.\ S., \& Strube, M. (2016).
Which coreference evaluation metric do you trust? A proposal for a link-based entity aware metric.
\textit{Proceedings of ACL 2016}, 632--642.

\bibitem{lappin1994algorithm}
Lappin, S., \& Leass, H.\ J. (1994).
An algorithm for pronominal anaphora resolution.
\textit{Computational Linguistics}, 20(4), 535--561.

\bibitem{lauscher2020zero}
Lauscher, A., Ravishankar, V., Vuli\'{c}, I., \& Glava\v{s}, G. (2020).
From zero to hero: On the limitations of zero-shot language transfer with multilingual transformers.
\textit{Proceedings of EMNLP 2020}, 4483--4499.

\bibitem{hedderich2021survey}
Hedderich, M.\ A., Lange, L., Adel, H., Sch\"{u}tze, H., \& Klakow, D. (2021).
A survey on recent approaches for natural language processing in low-resource scenarios.
\textit{Proceedings of NAACL 2021}, 2545--2568.

\bibitem{wolf2020transformers}
Wolf, T., Debut, L., Sanh, V., Chaumond, J., Delangue, C., Moi, A., \ldots\ \& Rush, A.\ M. (2020).
Transformers: State-of-the-art natural language processing.
\textit{Proceedings of EMNLP 2020: System Demonstrations}, 38--45.

\end{thebibliography}

% ========== PUBLICATIONS ==========
\chapter*{PUBLICATIONS FROM THE WORK}
\addcontentsline{toc}{chapter}{PUBLICATIONS FROM THE WORK}

\noindent No publications have resulted from this project to date. However, several aspects of the work represent potential contributions suitable for submission to appropriate venues. The documented failure analysis of unsupervised embedding-based coreference for Hindi, the manually constructed annotated dataset, and the empirical evaluation of the hybrid IndicBERTv2 pipeline together constitute a coherent set of findings that could form the basis of a paper targeting low-resource NLP workshops such as LoResMT or WILDRE (held alongside LREC-COLING), or student research tracks at venues including ACL, EMNLP, and IJCNLP. The primary steps needed before submission would be expanding the annotated dataset, implementing full CoNLL evaluation metrics, and conducting a more thorough comparison with existing Hindi coreference baselines.


% ========== APPENDICES ==========
\appendix

\chapter{CODE SNIPPETS}

\section{SentencePiece BPE Tokenizer Training (Phase 1)}

\begin{lstlisting}[language=Python, caption={SentencePiece BPE Tokenizer Training for Hindi}]
import sentencepiece as spm

spm.SentencePieceTrainer.train(
    input='hindi_wikipedia_corpus.txt',
    model_prefix='hindi_bpe_tokenizer',
    vocab_size=32000,
    character_coverage=0.9995,
    model_type='bpe',
    pad_id=0,
    unk_id=1,
    bos_id=2,
    eos_id=3,
    user_defined_symbols=['[MASK]', '[CLS]', '[SEP]']
)
\end{lstlisting}

\section{Scratch BERT Configuration (Phase 1)}

\begin{lstlisting}[language=Python, caption={Compact BERT Configuration for Scratch Training}]
from transformers import BertConfig, BertForMaskedLM

config = BertConfig(
    vocab_size=32000,
    hidden_size=256,
    num_hidden_layers=4,
    num_attention_heads=4,
    intermediate_size=512,
    max_position_embeddings=128,
    hidden_dropout_prob=0.1,
    attention_probs_dropout_prob=0.1
)

model = BertForMaskedLM(config)
print(f"Parameters: {model.num_parameters():,}")
\end{lstlisting}

\section{IndicBERTv2 Mention Representation (Phase 2)}

\begin{lstlisting}[language=Python, caption={Extracting Mention Embeddings from IndicBERTv2}]
from transformers import AutoTokenizer, AutoModel
import torch

tokenizer = AutoTokenizer.from_pretrained(
    "ai4bharat/indic-bert"
)
model = AutoModel.from_pretrained(
    "ai4bharat/indic-bert"
)
model.eval()

def get_mention_embedding(document, start, end):
    """Mean-pool last hidden states over mention span tokens."""
    inputs = tokenizer(
        document, return_tensors="pt",
        truncation=True, max_length=512
    )
    with torch.no_grad():
        outputs = model(**inputs)
    hidden = outputs.last_hidden_state[0]  # (seq_len, hidden)
    mention_hidden = hidden[start:end+1]
    return mention_hidden.mean(dim=0).numpy()
\end{lstlisting}

\section{Logistic Regression Pairwise Classifier (Phase 2)}

\begin{lstlisting}[language=Python, caption={Training the Pairwise Coreference Classifier}]
from sklearn.linear_model import LogisticRegression
from sklearn.preprocessing import StandardScaler
from sklearn.model_selection import StratifiedKFold
from sklearn.metrics import classification_report
import numpy as np

# X: feature matrix (n_pairs x n_features)
# y: labels (1 = coreferent, 0 = non-coreferent)

scaler = StandardScaler()
X_scaled = scaler.fit_transform(X)

clf = LogisticRegression(
    class_weight='balanced',
    solver='lbfgs',
    max_iter=1000,
    C=1.0
)

skf = StratifiedKFold(n_splits=5, shuffle=True, random_state=42)
all_preds, all_true = [], []

for train_idx, val_idx in skf.split(X_scaled, y):
    clf.fit(X_scaled[train_idx], y[train_idx])
    preds = clf.predict(X_scaled[val_idx])
    all_preds.extend(preds)
    all_true.extend(y[val_idx])

print(classification_report(all_true, all_preds,
      target_names=['Non-coref', 'Coreferent']))
\end{lstlisting}

% ========== APPENDIX B: DATASET ANNOTATION SCHEMA ==========
\chapter{DATASET ANNOTATION SCHEMA}

\section{Annotation Format}

Each document in the manually created dataset is stored as a JSON object following the schema shown below. The format was designed to be self-contained: every document carries its raw text, the full list of annotated mentions with their character offsets and types, and the cluster groupings as a separate dictionary indexed by cluster identifier.

\begin{lstlisting}[language=Python, caption={Dataset Annotation Schema (JSON)}]
{
  "doc_id": "doc_001",
  "text": "Ram school gaya. Vah khush tha...",
  "mentions": [
    {
      "mention_id": "m1",
      "start": 0,
      "end": 2,
      "text": "Ram",
      "mention_type": "proper_noun",
      "cluster_id": "C1"
    },
    {
      "mention_id": "m2",
      "start": 18,
      "end": 20,
      "text": "Vah",
      "mention_type": "pronoun",
      "cluster_id": "C1"
    }
  ],
  "clusters": {
    "C1": ["m1", "m2"]
  }
}
\end{lstlisting}

\section{Mention Type Taxonomy}

Table~\ref{tab:mention_types} describes the four mention types used during annotation. The taxonomy was chosen to cover the main categories of coreferring expressions found in Hindi text while remaining simple enough to apply consistently without extensive linguistic training.

\begin{table}[H]
\centering
\caption{Mention Type Taxonomy}
\label{tab:mention_types}
\begin{tabular}{|l|p{9cm}|}
\hline
\textbf{Type} & \textbf{Description} \\
\hline
\texttt{proper\_noun} & Named entity referring to a specific person, place, or organisation \\
\texttt{common\_noun} & Definite or indefinite common noun phrase functioning as an entity reference \\
\texttt{pronoun} & Personal, demonstrative, or relative pronoun \\
\texttt{nominal\_phrase} & Descriptive noun phrase that refers to an entity by role or attribute rather than name \\
\hline
\end{tabular}
\end{table}

\end{document}
